\section{The switched sum}\label{sec:l3}

We estimate the single term $-S_{\mathrm{sw}}/W$.
The normalization is \eqref{eq:not-normalization}.
Fix $\sigma$ as in \eqref{eq:not-sigma-choice}.
Put
\begin{equation}\label{eq:l3-constants}
\begin{gathered}
 \eta_L=10^{-4},\quad\varepsilon=10^{-4},\qquad
 H_{\max}=\frac{153}{256},\qquad
 \ell_0=\frac{264897}{2060000},\qquad \ell_1=\frac{109}{1000},\\
 N_{\mathrm{cut}}=N(1-L^{-333}),\qquad
 \mathcal P_N=\{p\in\mathbb P:p\nmid2N\}.
\end{gathered}
\end{equation}
All limits in this section fix $\sigma$ first.

\subsection{A common nonnegative source}

We use the prime-coordinate cells of Paper~II
\cite[(5.29a)--(5.31)]{PaperII}.
Their rules are
\begin{equation}\label{eq:l3-descriptors}
\begin{gathered}
 I_1=[\alpha,\beta),\quad I_2=[\beta,\lambda),\quad
 I_3=[\lambda,1/3),\quad I_4=[1/3,1],\\
 C_{\mathrm{L3}}=\frac{57}{64}-u-\frac{11\alpha}{5}
       =\frac{6408241}{9192000},\qquad
 v_\beta=\frac{2(C_{\mathrm{L3}}-\beta)}5,\\
 k_1+k_2+k_3+k_4=k,\quad 5\le k\le11,\quad
 k_1\ge m_1\ge\cdots\ge m_{k_2}\ge0,\\
 v_j\le v_\beta\ \Longrightarrow\ m_j=k_1,\\
 v_j>v_\beta\ \Longrightarrow\
 x_{m_j}<C_{\mathrm{L3}}-\frac52v_j\le x_{m_j+1},\\
 w_C=3+1_{k_1=0}-2k_1-2k_2-k_3+
            \binom{k_1}{2}+\sum_{j=1}^{k_2}m_j.
\end{gathered}
\end{equation}
The $x_i$ and $v_j$ increase in $I_1$ and $I_2$.
The outer sentinels are infinite.
Let $\mathcal C_+$ contain the cells with $w_C>0$.
For such a cell, put
\begin{equation}\label{eq:l3-source}
\begin{aligned}
 a_C(e)&=1_{\substack{e=p_1\cdots p_k,\ N^\alpha\le p_1<\cdots<p_k\\
             (\log_Np_1,\ldots,\log_Np_k)\in C}},\\
 A(e)&=\sum_{C\in\mathcal C_+}w_Ca_C(e)
                     1_{\varepsilon N<e<N_{\mathrm{cut}}},
                     &0\le A(e)&\le36,\\
 X_A&=\sum_eA(e),&Y_A(q)&=\sum_{(e,q)=1}A(e),\\
 C_A(q;N)&=\sum_eA(e)1_{q\mid N-e}.&&
\end{aligned}
\end{equation}
Since $\sum_jm_j\le k_1k_2$, we have
\[
 w_C\le3+1_{k_1=0}-2k_1-2k_2-k_3
             +\frac{k_1(k_1-1)}2+k_1k_2\le36.
\]
For $k_1\ge2$, maximize at $k_3=0$, $k_2=11-k_1$;
the resulting concave quadratic has maximum $36$ at $k_1=10,11$.
The cases $k_1=0,1$ give at most $4,1$.
The strict cells are disjoint, so this proves the bound for $A$.
Every prime is permitted in $A$.
In particular, $A$ has no factor $1_{(e,2N)=1}$.
Keep the selector $1_{\Omega(e)\ge3}$ in the switched source.
The required pointwise comparison is
\begin{equation}\label{eq:l3-domination}
 1_{\Omega(e)\ge3}(W_{13}(e))_+\le A(e)
 \quad\left(\begin{gathered}
 \mu_{\mathrm{Mob}}(e)^2=1,\quad P^-(e)\ge N^\alpha,\\
 \varepsilon N<e<N_{\mathrm{cut}}
 \end{gathered}\right).
\end{equation}
Here $\mu_{\mathrm{Mob}}$ denotes the M\"obius function.
Paper~II gives zero positive part when $\Omega(e)=3,4$
\cite[(1.7), (4.9), (5.29a)--(5.31)]{PaperII}.
\paragraph{Proof of \eqref{eq:l3-domination}.}
Let $e$ satisfy the side conditions and write $e=p_1\cdots p_k$,
$p_1<\cdots<p_k$, $x_i=\log_Np_i$.
If $(e,N)>1$, every occurrence in \eqref{eq:source-row-occurrences}
vanishes, because it requires $(d,N)=1$ and $(n/d,Nd)=1$
\cite[(1.7)]{PaperII}. Then $W_{13}(e)=0$. Assume $(e,N)=1$.
Since $N^{k\alpha}\le e<N$ and $12\alpha=1200/1149>1$, we have $k\le11$.
If $k\le2$, the selector vanishes. If $k\in\{3,4\}$, Paper~II, (4.9),
gives $W_{13}(e)\le0$. Its hypotheses are exactly squarefree $e$,
$(e,N)=1$, $\varepsilon N<e<N$ and $P^-(e)\ge N^\alpha$
\cite[(4.9)]{PaperII}.
Let $5\le k\le11$. The table after Paper~II, (5.29a), shows that the
three negative pair counts vanish, so
\[
 W_{13}(e)=3+1_{k_1=0}-2k_1-2k_2-k_3+\binom{k_1}2+\sum_{j=1}^{k_2}m_B(v_j),
\]
where $m_B(v)$ counts first-band coordinates below $B_{\mathrm{O1}}(v)$
\cite[(5.29a), (5.30)]{PaperII}. The bins determine $k_1,\ldots,k_4$.
If $v_j\le v_\beta$, then $B_{\mathrm{O1}}(v_j)=\beta$ exceeds every
first-band coordinate, so $m_B(v_j)=k_1$. If $v_j>v_\beta$, then
$B_{\mathrm{O1}}(v_j)=C_{\mathrm{L3}}-\frac52v_j$, and $m_B(v_j)$ is the unique $m_j$
with $x_{m_j}<C_{\mathrm{L3}}-\frac52v_j\le x_{m_j+1}$ (infinite sentinels).
It is nonincreasing in $j$. Thus $x$ lies in exactly one cell $C$ of
\eqref{eq:l3-descriptors}, and $W_{13}(e)=w_C$
\cite[(5.31)]{PaperII}. If $w_C\le0$, the left side of
\eqref{eq:l3-domination} is zero. Otherwise $C\in\mathcal C_+$ and
$a_C(e)=1$, so the left side equals $w_C=A(e)$. This proves
\eqref{eq:l3-domination}. The same argument, without the upper cut, gives
\begin{equation}\label{eq:l3-domination-nocut}
 1_{\Omega(e)\ge3}(W_{13}(e))_+\le A^0(e):=\sum_{C\in\mathcal C_+}w_Ca_C(e)\le36
\end{equation}
for squarefree $e$ with $P^-(e)\ge N^\alpha$ and $\varepsilon N<e<N$.

Put
\begin{equation}\label{eq:l3-sieved-source}
\begin{gathered}
 b_A(n)=A(N-n),\qquad
 P_N(z)=\prod_{\substack{p<z\\p\in\mathcal P_N}}p,\\
 S_v(z)=\sum_{\substack{n\ge1\\v\mid n}}b_A(n)1_{(n,P_N(z))=1},
 \qquad S_A=S_1(N^{\mu_*}),\\
 E_{\mathrm{src}}=(S_{\mathrm{sw}}-S_A)_+,\qquad
 Z_{\mathrm{src}}=(S_A-S_{\mathrm{sw}})_+.
\end{gathered}
\end{equation}
The definitions give the finite identity
\begin{equation}\label{eq:l3-source-id}
 S_{\mathrm{sw}}=S_A+E_{\mathrm{src}}-Z_{\mathrm{src}},
 \qquad E_{\mathrm{src}},Z_{\mathrm{src}}\ge0.
\end{equation}
Under \eqref{eq:l3-domination}, the required bound is
\begin{equation}\label{eq:l3-source-error}
 E_{\mathrm{src}}\le910N^{1-\alpha}
            +36(NL^{-333}+1)+36N^{\mu_*}.
\end{equation}
There are at most $2N^{1-\alpha}$ rough squareful integers.
Each costs at most $455$, by the pointwise weight bound of
Paper~II, proof of Lemma~1.1, immediately after (1.11)
\cite[Lemma~1.1, proof after (1.11)]{PaperII}.
The last two terms cover the top strip and small primes.
\paragraph{Proof of \eqref{eq:l3-source-error}.}
Write $e=N-p$ for $p\in\mathbb P$, $p<(1-\varepsilon)N$.
Every prime $p\ge N^{\mu_*}$ has $(p,P_N(N^{\mu_*}))=1$, and $A\ge0$.
Hence $S_A\ge\sum_{p\ge N^{\mu_*}}A(N-p)$, and
\[
 S_{\mathrm{sw}}-S_A\le\sum_{p\ge N^{\mu_*}}\bigl(s_{\mathrm{sw}}(e)-A(e)\bigr)_+
                +\sum_{p<N^{\mu_*}}s_{\mathrm{sw}}(e).
\]
Here $s_{\mathrm{sw}}(e)=0$ unless $(e,N)=1$ and $P^-(e)\ge N^\alpha$,
and always $0\le s_{\mathrm{sw}}(e)\le|W_{13}(e)|\le455$
\cite[Lemma~1.1, proof after (1.11)]{PaperII}.
Split the summands into three classes; overlaps only enlarge the bound.
(i) Non-squarefree $e<N$ with $P^-(e)\ge N^\alpha$ number at most
$\sum_{\ell\ge N^\alpha}N/\ell^2\le2N^{1-\alpha}$; they cost at most
$455\cdot2N^{1-\alpha}=910N^{1-\alpha}$.
(ii) Squarefree $e$ with $N_{\mathrm{cut}}\le e<N$ number at most
$NL^{-333}+1$; by \eqref{eq:l3-domination-nocut} they cost at most
$36(NL^{-333}+1)$.
(iii) The primes $p<N^{\mu_*}$ with squarefree $e$ cost at most $36N^{\mu_*}$,
again by \eqref{eq:l3-domination-nocut}.
For the remaining squarefree $e$ with $\varepsilon N<e<N_{\mathrm{cut}}$
and $p\ge N^{\mu_*}$, \eqref{eq:l3-domination} gives a zero summand.
This proves \eqref{eq:l3-source-error}. Since $W\gg N/L^2$,
$E_{\mathrm{src}}/W\ll L^{-331}$.

\subsection{Four rows and the lower-sieve cutoff}

Use the clipped intervals
$[N^t,N^{\rho t})\cap[N^u,N^{\mu_*})$.
Let $\mathcal A_j$ comprise the ordered prefixes
\[
 N^u\le a_j<\cdots<a_1<N^{\mu_*},\qquad a_i\in\mathcal P_N.
\]
Write $v=a_1\cdots a_j$.
The empty prefix has $v=1$.
Define
\begin{equation}\label{eq:l3-buchstab}
\begin{gathered}
 T_0=S_1(N^u),\quad
 T_1=\sum_{\mathcal A_1}S_{a_1}(N^u),\quad
 T_2=\sum_{\mathcal A_2}S_{a_1a_2}(N^u),\\
 T_3^{\mathrm{dyn}}=\sum_{\mathcal A_3}S_{a_1a_2a_3}(a_3),
 \qquad S_A=T_0-T_1+T_2-T_3^{\mathrm{dyn}}.
\end{gathered}
\end{equation}
This is Buchstab's finite identity
\cite[(4.18)]{PaperII}.
Its coefficient is
$1-k+\binom{k}{2}-\binom{k-1}{2}=0$ for $k\ge1$.
At $k=0$ it equals $1$.

At depth three, let $\Theta_c$ be the infimum
of $\widetilde H$ on the closed clipped parent. The underlying prime
cells stay half-open. The closure is used only for this infimum.
Use the whole anatomy function before choosing a branch
\cite[(2.17d)--(2.17e), (4.19)]{PaperII}.
Set
\begin{equation}\label{eq:l3-row3-level}
 \ell_c=\Theta_c-\eta_L-t_1-t_2-t_3,\quad
 D_c=N^{\ell_c},\quad z_c=N^{t_3},\quad
 \chi_c=1_{\ell_c\ge2\rho t_3}.
\end{equation}
Retained prefixes have $\chi_c=1$.
For these prefixes define
\begin{equation}\label{eq:l3-peel}
\begin{aligned}
 B_3&=\sum_{\mathcal A_3:\chi_c=1}S_v(z_c),&
 D_-&=\sum_{\mathcal A_3:\chi_c=0}S_v(a_3),\\
 D_+&=\sum_{\mathcal A_3:\chi_c=1}\{S_v(z_c)-S_v(a_3)\},&
 P_3&=\sum_{\mathcal A_3:\chi_c=1}
       \sum_{\substack{r\in\mathcal P_N\\z_c\le r<a_3}}S_{vr}(z_c),\\
 Z_{\mathrm{peel}}&=D_-+P_3-D_+\ge0.&&
\end{aligned}
\end{equation}
The union bound gives $D_+\le P_3$.
Therefore
\begin{equation}\label{eq:l3-buchstab-peel}
 S_A=T_0-T_1+T_2-B_3+P_3-Z_{\mathrm{peel}}.
\end{equation}
The favorable term enters exactly once
\cite[(4.45)]{PaperII}.

\subsection{Capping the first three levels}

At depths $j=0,1,2$, replace the relative level by
\begin{equation}\label{eq:l3-cap}
 T_c=\sum_{i=1}^jt_i,\quad
 \ell_c=\frac12-\eta_L-\rho T_c,\quad
 D_c=N^{\ell_c},\quad z_{j,c}=N^u.
\end{equation}
Here $T_c=0$ at depth zero.
Depth three keeps \eqref{eq:l3-row3-level}, with $z_{3,c}=z_c$.
The elementary bounds are
\begin{equation}\label{eq:l3-cap-margins}
\begin{gathered}
 \ell_c\ge\frac12-\eta_L-j\rho\mu_*\ge\ell_1,
 \qquad s_c=\ell_c/u\ge54.6375-10^{-87}>30,\\
 \text{continuum reference values for depths }0,1,2\text{ (the first three rows)}
 \text{ at }(\eta_L,\rho)=(0,1):
       250,\ \frac{4875}{32},\ \frac{875}{16},\\
 q<N^{\rho T_c}D_c^{\rho+\sigma}
   \le N^{1/2-\eta_L+(\sigma+\tau)/2}
   <N^{1/2-9\cdot10^{-5}}.
\end{gathered}
\end{equation}
The modulus includes the smooth divisor and every tail prime.
The exponent $\rho+\sigma$ is kept.
The cap changes no row-three level.

The lower floor changes the repeated-prime argument.
Put
\begin{equation}\label{eq:l3-rough-floor}
 u_1=N^{\ell_1\sigma^2},\qquad
 M_\sigma=\left\lceil\frac3{5\ell_0\sigma^2}\right\rceil,
 \qquad C_{\mathrm{L3},\sigma}=M_\sigma!.
\end{equation}
\paragraph{(P2.8b) Rough-coordinate and repeated-prime bounds.}
Assume the support and rough-coordinate hypotheses stated below.
These are the capped version of Paper~II's Lemmas B.1--B.4.
The coefficients and masks are defined in \eqref{eq:l3-masks}.
Every rough coordinate is at least $u_c=D_c^{\sigma^2}\ge u_1$.
On an active block, count these coordinates by $m_c$.
The capped modulus bound gives
\begin{equation}\label{eq:l3-rough-count}
 u_1^{m_c}\le q<N^{1/2},\qquad
 m_c<\frac{500}{109\sigma^2}
       \le\frac3{5\ell_0\sigma^2},\qquad
 \frac{m_c!}{M_c}\le C_{\mathrm{L3},\sigma}.
\end{equation}
The prefix orbit size is $M_c\in\{1,2,6\}$.
The small-prime part of $q$ fixes the smooth divisor.
The rough multiset fills at most $m_c!$ labelled positions.
Thus
\begin{equation}\label{eq:l3-cap-coefficient-bound}
 \frac{|C_{\mathfrak b}(q)|}{M_c}\le C_{\mathrm{L3},\sigma}.
\end{equation}
Define $\mathcal H_1=\{q<N^{1/2}:p^2\mid q\text{ for some }p\ge u_1\}$.
The forbidden and allowed-repeat masks satisfy
\begin{equation}\label{eq:l3-u1-masks}
 |\Xi_{U,\mathfrak b}(q)|+|\Xi_{A,\mathfrak b}(q)|
                \le C_{\mathrm{L3},\sigma}1_{\mathcal H_1}(q).
\end{equation}
Writing $q=vp^2$ and overcounting gives
\begin{equation}\label{eq:l3-u1-sums}
 \sum_{q\in\mathcal H_1}\frac1q\le\frac{2(1+L)}{u_1},
 \qquad \#\mathcal H_1\le\frac{2N^{1/2}}{u_1}.
\end{equation}
The tested measure has $0\le A\le36$ and support below $N$.
For $N\ge N_{\mathrm{dom}}$, $q/\varphi(q)\le1+L$ throughout $q<N^{1/2}$.
Consequently
\begin{equation}\label{eq:l3-u1-test}
 |\Delta_P(q)|,\ |d_A(q)|
 \le36\left(\frac Nq+1+\frac N{\varphi(q)}\right)
 \le36\left(\frac{N(2+L)}q+1\right).
\end{equation}
Per block, these bounds give
\begin{equation}\label{eq:l3-u1-block}
 \sum_q\bigl(|\Xi_U(q)\Delta_P(q)|+|\Xi_A(q)d_A(q)|\bigr)
       \le\frac{216C_{\mathrm{L3},\sigma} N(1+L)^2}{u_1}.
\end{equation}
There are finitely many blocks at fixed $\sigma$.
Division by $W\gg N/L^2$ proves the conditional estimate
\begin{equation}\label{eq:l3-repeated-cap}
 |E_{\mathrm{sq}}|/W=O_\sigma(N^{-\ell_1\sigma^2}L^4).
\end{equation}
These are Paper~II, Lemmas B.1--B.4 and (B.3)--(B.7),
with $u_0$ replaced by $u_1$ on rows $0,1,2$
\cite[Appendix~B.1, Lemmas B.1--B.4, (B.3)--(B.7)]{PaperII}.
Row three keeps $u_0=N^{\ell_0\sigma^2}$, since its level floor
in \eqref{eq:l3-row3-global-margins} exceeds $\ell_0$.
\paragraph{Verification for the capped rows.}
Fix a block in row $j\le2$. Prefix primes lie in
$[N^{t_i},N^{\rho t_i})\cap[N^u,N^{\mu_*})$, so their product is below
$N^{\rho T_c}$. The tail tuple satisfies $\prod_iD_i<D_c$ by its last
cubic constraint, and the tail primes are below $D_i^\rho$; also
$d<D_c^\sigma$. Thus
\[
 \log_Nq<\rho T_c+(\rho+\sigma)\ell_c
 =(\rho+\sigma)(\tfrac12-\eta_L)-\rho(\rho+\sigma-1)T_c
 \le(\rho+\sigma)(\tfrac12-\eta_L)<\tfrac12-9\cdot10^{-5}.
\]
Every prime of $q$ lies in $\mathcal P_N$ and $d\mid P_N(u_c)$, so
$(q,2N)=1$. Tail primes are at least $D_i\ge D_c^{\sigma^2}=u_c$.
Prefix primes are at least $N^u>u_c$, because $\ell_c\sigma^2<u$.
Also $u_c=N^{\ell_c\sigma^2}\ge N^{\ell_1\sigma^2}=u_1$.
For an active block choose a nonzero realization. Then
$u_1^{m_c}\le q<N^{1/2}$, so $m_c<1/(2\ell_1\sigma^2)=500/(109\sigma^2)$,
and $1/(2\ell_1)\le3/(5\ell_0)$. This is \eqref{eq:l3-rough-count}.
The primes of $q$ below $u_c$ determine $d$. The rough multiset fills
the $m_c$ labelled positions in at most $m_c!$ ways, and $|\lambda_d|\le1$.
This gives \eqref{eq:l3-cap-coefficient-bound}. It also gives
\eqref{eq:l3-u1-masks}. A forbidden tuple repeats a prime $\ge u_c$.
An allowed non-squarefree tuple has $p^2\mid q$ with $p\ge u_c$,
since $d$ is squarefree and coprime to the rough part. The grid count of
Paper~II, Lemma~B.1(4), gives finitely many blocks at fixed $\sigma$.
These are Paper~II, Lemmas~B.1--B.2, with $(u_0,N^{3/5})$ replaced by
$(u_1,N^{1/2})$ \cite[Appendix~B.1, Lemmas~B.1--B.2]{PaperII}.
Lemma~B.4 with $u_1$ and $Q=N^{1/2}$ gives \eqref{eq:l3-u1-sums}.
Lemma~B.3, with $A\le36<A_{\max}$, support in $[1,N)$ and
$q/\varphi(q)\le1+\log_2q\le1+L$, gives \eqref{eq:l3-u1-test}.
Then $\sum_{q\in\mathcal H_1}(N(2+L)/q+1)\le6N(1+L)^2/u_1$ gives
\eqref{eq:l3-u1-block}. Row three is Paper~II, Lemmas~B.1--B.4 as
printed: its modulus bound $q<N^{e_Q}<N^{3/5}$ and floor
$\ell_c>\ell_0$ are \eqref{eq:l3-row3-global-margins}. Hence
\eqref{eq:l3-repeated-cap} holds for all four rows.

\subsection{The signed sieve identity}

Put $\nu_j=(-1)^j$.
For all four rows, define
\begin{equation}\label{eq:l3-common-sifting-ratio}
 s_c=\frac{\ell_c}{\log_N z_{j,c}}
 =\begin{cases}\ell_c/u,&j=0,1,2,\\ \ell_c/t_3,&j=3.\end{cases}
\end{equation}
Let $K_j=F$ for even $j$, and $K_j=f$ for odd $j$.
For a prefix $v$, set
\begin{equation}\label{eq:l3-density}
\begin{gathered}
 \mathcal X_v=X_A/\varphi(v),\qquad
 \omega_v(p)=\frac p{p-1}1_{p\in\mathcal P_N},
 \qquad \omega_v(p^h)=0\quad(h\ge2),\\
 V_N(z)=\prod_{\substack{p<z\\p\in\mathcal P_N}}
                     \left(1-\frac1{p-1}\right),\qquad
 \Delta_0(q)=C_A(q;N)-\frac{X_A}{\varphi(q)}.
\end{gathered}
\end{equation}
All cutoff primes lie below the prefix primes.
Thus the density product on the sieve support is $V_N$.
The squarefree remainder is $\Delta_0(vd)$.
Apply Fact~\ref{fact:structured-sieve} with
$\epsilon_s=\sigma$, $\delta_s=\tau$, and $L_\omega=2$.
Choose a common nonnegative majorant
\begin{equation}\label{eq:l3-sieve-error}
 E_{s,c}:= C_E\{\sigma+
     \sigma^{-8}e^{K_{\mathrm M}+2}(\ell_cL)^{-1/3}\}.
\end{equation}
The signed sieve error satisfies \(|E_s|\le E_{s,c}\).
\paragraph{Verification of the sieve data.}
Fix a row, a cell and a prefix $v$, and sieve the finite multiset
$\{n:v\mid n\}$ with multiplicity $b_A(n)\ge0$ by $\mathcal P_N$ below
$z=z_{j,c}$. Every prefix prime is at least $z$: on rows $0,1,2$ it is at
least $N^u=z$, and on row three $a_3\ge N^{t_3}=z_c$. Hence $(d,v)=1$
for $d\mid P_N(z)$. In particular $\omega_v(d)=\prod_{p\mid d}p/(p-1)$
for these $d$; the primes $p\mid v$ are all $\ge z$ and are never sifted.
For squarefree such $d$,
\[
 |\mathcal A_d|=\sum_eA(e)1_{vd\mid N-e}=C_A(vd;N),\qquad
 \frac{\omega_v(d)}d\mathcal X_v=\frac{X_A}{\varphi(v)\varphi(d)}
                               =\frac{X_A}{\varphi(vd)},
\]
so $r(\mathcal A,d)=\Delta_0(vd)$. Each remainder modulus in
\eqref{eq:l3-remainders} is squarefree. Here $d$ uses primes below
$D_c^{\sigma^2}$. Tail primes are at least $D_c^{\sigma^2}$ and distinct
(equal labels by the mask, unequal grid intervals disjoint). Prefix
primes are at least $z$. Next, $0<\omega_v(p)=p/(p-1)<p$ for every
$p\in\mathcal P_N$ (all such $p$ are odd),
$\omega_v=0$ off $\mathcal P_N$, and $\omega_v(p^h)=0$ for $h\ge2$. Hence the
second condition in \eqref{eq:classical-density} holds with $L_\omega=2$.
For the first, $(1-\omega_v(p)/p)^{-1}=(1-1/(p-1))^{-1}$, and Paper~II,
(5.5), gives
\[
 \prod_{\substack{w\le p<z\\p\in\mathcal P_N}}\left(1-\frac1{p-1}\right)^{-1}
 \le\frac{\log z}{\log w}\left(1+\frac{K}{\log w}\right)\qquad(2\le w<z)
\]
with one absolute $K$, uniformly in $N$ \cite[(5.5)]{PaperII}.
Take $K_{\mathrm M}$ in \eqref{eq:l3-sieve-error} to be the maximum of
this $K$ and the constant of \eqref{eq:classical-mertens-density}. It
remains absolute and is fixed before $\sigma$. The product
$V(z)$ of Fact~\ref{fact:structured-sieve} is then $V_N(z)$.
The levels satisfy $z\le D_c$ on rows $0,2$ ($s_c>30$) and
$z\le D_c^{1/2}$ on rows $1,3$ ($s_c>30$, resp.\ $\chi_c=1$, i.e.\
$\ell_c\ge2\rho t_3$). Fact~\ref{fact:structured-sieve} with
$\epsilon_s=\sigma$, $\delta_s=\tau$, $D=D_c$ therefore applies to every
actual sequence, and \eqref{eq:l3-signs} follows with
$E_s\le E_{s,c}$, $\log D_c=\ell_cL$.

For each cell fix one coefficient construction.
Use it for every prefix realization in that cell.
Write its remainders as
\begin{equation}\label{eq:l3-remainders}
\begin{aligned}
 R_{1,v}^{\nu}&=\sum_{\substack{d<D_c^\sigma\\d\mid P_N(D_c^{\sigma^2})}}
                         \phi_d^\nu\Delta_0(vd),\\
 R_v^\nu&=\sum_{\mathbf D\in\mathcal D^\nu(D_c)}
   \sum_{\substack{d<D_c^\sigma\\d\mid P_N(D_c^{\sigma^2})}}
   \lambda_d^\nu(\mathbf D)
   \sum_{\substack{D_i\le p_i<\min(D_i^\rho,z_{j,c})\\
                p_i\in\mathcal P_N,\ \text{equal-label entries distinct}}}
                         \Delta_0(vd\textstyle\prod_i p_i).
\end{aligned}
\end{equation}
The coefficients satisfy $|\phi_d^\nu|,|\lambda_d^\nu(\mathbf D)|\le1$
\cite[Theorem~2.6, (2.29)--(2.33b)]{PaperI}.
Define the finite quantities
\begin{equation}\label{eq:l3-four-rows}
\begin{aligned}
 M_j&=\sum_c\sum_{v\in c}\frac{X_A}{\varphi(v)}
                   V_N(z_{j,c})K_j(s_c),\\
 E_j&=\sum_c\sum_{v\in c}\frac{X_A}{\varphi(v)}
                   V_N(z_{j,c})E_{s,c},\qquad
 R_j=\sum_c\sum_{v\in c}(R_{1,v}^{\nu_j}+R_v^{\nu_j}),\\
 Q_j&=M_j+\nu_j E_j+\nu_j R_j,\qquad
 (S_0,S_1,S_2,S_3)=(T_0,T_1,T_2,B_3),\\
 Z_{\mathrm{AA}}&=\sum_{j=0}^3\nu_j(Q_j-S_j),\qquad
 M_{\mathrm{prin}}=\sum_{j=0}^3\nu_jM_j,\\
 M_{\mathrm{err}}&=\sum_{j=0}^3E_j,\qquad
 R_{\mathrm{adv}}=\sum_{j=0}^3R_j.
\end{aligned}
\end{equation}
The sums keep strict prefix order.
Under the stated sieve hypotheses,
\begin{equation}\label{eq:l3-signs}
 Q_0\ge S_0,\quad Q_1\le S_1,\quad
 Q_2\ge S_2,\quad Q_3\le S_3,\quad Z_{\mathrm{AA}}\ge0.
\end{equation}
Multiplying the signs gives
\begin{equation}\label{eq:l3-aa-id}
 \sum_{j=0}^3\nu_jS_j
    =M_{\mathrm{prin}}+M_{\mathrm{err}}+R_{\mathrm{adv}}-Z_{\mathrm{AA}}.
\end{equation}
Every remainder in $R_{\mathrm{adv}}$ has sign $+1$.

\subsection{The natural coprime comparator}

A block $\mathfrak b=(j,c,f,\mathbf D)$ fixes a row,
a cell, and either remainder type $f$.
The tail vector is empty for $R_1$.
Prefix realizations stay inside the coefficient.
Let $K_P,K_T$ impose prefix and tail distinctness.
Let $C_{\mathfrak b}$ be the unrestricted labelled coefficient.
Define
\begin{equation}\label{eq:l3-masks}
\begin{aligned}
 C_{\mathfrak b}(q)&=\sum_{\mathrm{full\ labelled\ tuples}}
          \widetilde\lambda_d1_{q=d\prod p_i},\\
 O_{\mathfrak b}(q)&=\frac1{M_c}
      \sum_{\mathrm{full\ labelled\ tuples}}
          \widetilde\lambda_dK_PK_T1_{q=d\prod p_i},\\
 \Xi_{U,\mathfrak b}(q)&=\frac1{M_c}
      \sum_{\mathrm{full\ labelled\ tuples}}
          \widetilde\lambda_d(1-K_PK_T)1_{q=d\prod p_i},\\
 O_{\mathfrak b}&=C_{\mathfrak b}/M_c-\Xi_{U,\mathfrak b},\qquad
 \Xi_{A,\mathfrak b}=O_{\mathfrak b}1_{q\ \mathrm{not\ squarefree}}.
\end{aligned}
\end{equation}
The two comparator discrepancies are
\begin{equation}\label{eq:l3-comparator}
 \Delta_P(q)=C_A(q;N)-\frac{Y_A(q)}{\varphi(q)},\qquad
 d_A(q)=\Delta_0(q)-\Delta_P(q)
             =-\frac{X_A-Y_A(q)}{\varphi(q)}.
\end{equation}
For $\langle f,g\rangle=\sum_qf(q)g(q)$, put
\begin{equation}\label{eq:l3-adv-split}
\begin{aligned}
 E_{\mathrm{disp},f}&=\sum_{\mathfrak b:\mathrm{type}=f}
                \frac{\langle C_{\mathfrak b},\Delta_P\rangle}{M_c}
                     \quad(f=1,2),\\
 E_{\mathrm{sq}}&=\sum_{\mathfrak b}
       \{-\langle\Xi_{U,\mathfrak b},\Delta_P\rangle
                         +\langle\Xi_{A,\mathfrak b},d_A\rangle\},\\
 E_{\mathrm{nc}}&=\sum_{\mathfrak b}\sum_{q\ \mathrm{squarefree}}
                              O_{\mathfrak b}(q)d_A(q),\\
 R_{\mathrm{adv}}&=E_{\mathrm{disp},1}+E_{\mathrm{disp},2}
                                      +E_{\mathrm{sq}}+E_{\mathrm{nc}}.
\end{aligned}
\end{equation}
This is a finite coefficient identity
\cite[(4.22)--(4.26), Appendix~B.1--B.4]{PaperII}.
All moduli satisfy $(q,2N)=1$.
That restriction does not puncture $A$.

\subsection{The full-support bilinear application}

Split each prime coordinate into relative-width $L^{-334}$ intervals.
An interior box has constant order and descriptor predicates.
For $B=\prod_i[A_i,B_i)$, put
\begin{equation}\label{eq:l3-direct-prime}
\begin{gathered}
 T=\lceil A_1\rceil-1,\quad U_J=\lceil B_1\rceil-1,\quad
 I_T=(T,U_J],\quad X=\prod_iA_i,\\
 a_n=1_{\mathbb P}(n)1_{I_T}(n),\qquad
 b_m=\sum_{\substack{p_2<\cdots<p_k\\p_i\in[A_i,B_i)}}1_{m=p_2\cdots p_k},\\
 b_m\in\{0,1\},\qquad
 N^\alpha/2\le T<N^{1/5},\qquad \varepsilon N\le X\le N.
\end{gathered}
\end{equation}
The box contributes exactly $w_Ca_nb_m$.
For $T>2^{50}$ and $U_J\le2T$,
\begin{equation}\label{eq:l3-prime-option}
 a_n=1_{P^-(n)\ge T^{51/100}}1_{I_T}(n).
\end{equation}
A composite $n\le2T$ has a smaller prime divisor.
This is Paper~I, (3.8), option 1
\cite[(3.8)--(3.9)]{PaperI}.
Its logarithmic replacement error is zero.

Define the full-support discrepancy
\begin{equation}\label{eq:l3-box-discrepancy}
\begin{aligned}
 \Delta_B(q,N)&=\sum_{nm\equiv N\ (q)}a_nb_m
           -\frac1{\varphi(q)}\sum_{(nm,q)=1}a_nb_m,\\
 \mathcal S_B(\mathfrak b)&=\sum_qC_{\mathfrak b}(q)\Delta_B(q,N).
\end{aligned}
\end{equation}
There is no condition $nm\le X$ in these sums.
The exact interior decomposition is
\begin{equation}\label{eq:l3-interior-split}
\begin{aligned}
 D_{\mathrm{int}}^{(j)}
 &=\sum_{\mathfrak b:\mathrm{row}=j}\frac1{M_c}
           \sum_C\sum_{B\subset C\ \mathrm{interior}}w_C\mathcal S_B(\mathfrak b),\\
 E_{\mathrm{disp},1}+E_{\mathrm{disp},2}
 &=\sum_{j=0}^3D_{\mathrm{int}}^{(j)}+D_{\mathrm{bdry}}.
\end{aligned}
\end{equation}

For rows $0,1,2$, apply Fact~\ref{fact:bv-full-support}.
Its Siegel--Walfisz sequence must be the selected prime sequence.
The dictionary with Fact~\ref{fact:bv-full-support} is
$b_n^{\mathrm{F2}}:=a_n$, $a_m^{\mathrm{F2}}:=b_m$, and $Y:=T$.
Thus the prime sequence carries the Siegel--Walfisz hypothesis.
We need Paper~I, (3.25), in the form
\begin{equation}\label{eq:l3-prime-sw}
 \sum_{\substack{n\in I\\n\equiv l\ (k)\\(n,d)=1}}1_{\mathbb P}(n)
 -\frac1{\varphi(k)}\sum_{\substack{n\in I\\(n,dk)=1}}1_{\mathbb P}(n)
 \ll_{A_{\mathrm{SW}}}T\tau(d)L^{-A_{\mathrm{SW}}},\qquad(l,k)=1,
 \quad I\subset I_T.
\end{equation}
The normalized hypothesis also needs
\begin{equation}\label{eq:l3-prime-norm}
 \|a\|_2T^{1/2}\gg T L^{-168}.
\end{equation}
Use Fact~\ref{fact:prime-estimates}, specifically \eqref{eq:classical-sw}.
Its endpoint estimate gives, for every fixed $A>335$,
\begin{equation}\label{eq:l3-endpoint-prime-count}
 \pi(U_J)-\pi(T)=\operatorname{Li}(U_J)-\operatorname{Li}(T)
                         +O_A(TL^{-A}).
\end{equation}
Full mesh intervals satisfy $U_J-T\gg TL^{-334}$.
Thus $\#(I_T\cap\mathbb P)\gg TL^{-335}$ uniformly.
This implies \eqref{eq:l3-prime-norm}.
The $k=1$ discrepancy in \eqref{eq:l3-prime-sw} is identically zero.
The counting input is the endpoint estimate just displayed.
\paragraph{Proof of \eqref{eq:l3-prime-sw} and \eqref{eq:l3-prime-norm}.}
Let $N^\alpha/2\le T<N^{1/5}$, $U_J\le2T$, $(l,k)=1$, and let
$I=(y_1,y_2]\subset I_T$ be an arbitrary subinterval, so
$T\le y_1\le y_2\le U_J$. (Paper~I, (3.25), is quantified over all
such $I$ \cite[(3.25)]{PaperI}.)
The condition $(n,d)=1$ removes at most $\omega(d)$ primes from each sum.
If $k\le L^{A+2}$, apply \eqref{eq:classical-sw} at $y=y_2$ and $y=y_1$,
with $C=A+3$. Since $T\le y\le2T$, $\log y\asymp L$ and the error is
$O_A(TL^{-A})$ at each endpoint. Hence
$\pi(y_2;k,l)-\pi(y_1;k,l)=(\operatorname{Li}(y_2)-\operatorname{Li}(y_1))/\varphi(k)+O_A(TL^{-A})$.
The same holds for the unrestricted count divided by $\varphi(k)$,
up to the $\omega(k)$ primes dividing $k$.
If $k>L^{A+2}$, since $y_2-y_1\le T$ both terms are at most
$T/k+1+T/\varphi(k)\ll TL^{-A}$,
using $\varphi(k)\gg k/\log\log k$. This proves \eqref{eq:l3-prime-sw}
with $\tau(d)$, ineffectively, uniformly in $I\subset I_T$. For the norm,
$U_J-T\ge A_1L^{-334}-1\ge TL^{-334}/2$ on a full mesh interval, and
\eqref{eq:l3-endpoint-prime-count} with $A=336$ gives
$\#(I_T\cap\mathbb P)\ge TL^{-335}/4$ eventually. Hence
$\|a\|_2T^{1/2}\ge TL^{-335/2}/2\ge TL^{-168}$. Both bounds are uniform
over all interior boxes, since only $T$ enters.

Assume \eqref{eq:l3-prime-sw} and \eqref{eq:l3-prime-norm} uniformly.
Taking $A_{\mathrm{SW}}=A_0+168$ verifies
\eqref{eq:bvfs-sw} with divisor exponent one.
Both factor lengths exceed $x_B^{1/12}$ eventually.
Split each of the $k-1$ complementary prime intervals at any dyadic
boundary. This gives at most $2^{k-1}$ pieces. Localize each resulting
product on $m\sim M_B$. Since $(1+L^{-334})^{k-1}<2$ eventually,
there are at most two product dyadic intervals per piece.
This bounded splitting changes no logarithmic exponent.
Set $x_B=TM_B\asymp N$.
There are at most $T$ selected integers in $(T,2T]$ and at most
$2M_B$ integers in a product dyadic interval. Unique factorization gives
$\|a\|_2\|b\|_2\le\sqrt{2x_B}$.
The required application is
\begin{equation}\label{eq:l3-bv-application}
\begin{gathered}
 \frac{|C_{\mathfrak b}(q)|}{M_c}\le C_{\mathrm{L3},\sigma},\qquad
 Q<N^{1/2-9\cdot10^{-5}}\le x_B^{1/2}(\log x_B)^{-B_{\mathrm{BV}}},\\
 \|a\|_2\|b\|_2\le\sqrt{2}\,x_B^{1/2},\qquad
 |\mathcal S_B(\mathfrak b)|/M_c\ll_\sigma NL^{-3692}.
\end{gathered}
\end{equation}
Use $A=3692$, $B=1$, $\epsilon=1/12$ and
$B_{\mathrm{BV}}=B_{\mathrm{BV}}(3692,1,1/12)$ in \eqref{eq:bvfs}.
The source has $O_\sigma(L^{3689})$ weighted boxes.
There are finitely many blocks at fixed $\sigma$.
Since $W\gg N/L^2$, summation gives
\begin{equation}\label{eq:l3-bv-total}
 \frac{|D_{\mathrm{int}}^{(0)}|+|D_{\mathrm{int}}^{(1)}|
          +|D_{\mathrm{int}}^{(2)}|}{W}\ll_\sigma L^{-1}.
\end{equation}
\paragraph{Verification of \eqref{eq:l3-bv-application}.}
The coefficient bound is \eqref{eq:l3-cap-coefficient-bound}, proved above,
and the modulus range is \eqref{eq:l3-cap-margins}. Every supported $q$ is
prime to $2N$, so $a=N$ is a reduced class and the maximum in
\eqref{eq:bvfs} covers it. The selected sequence $a_n$ is supported in
$(T,2T]$, and $b_m$, after at most two dyadic pieces, in $(M_B,2M_B]$.
For $x_B=TM_B$, $x_B\in[\varepsilon N/8,2N]$ and
$\log T/\log x_B\in[\alpha-o(1),1/5+o(1)]\subset[1/12,11/12]$, because
$\alpha-1/12=17/4596>0$. The hypothesis \eqref{eq:bvfs-sw} holds with
$B=1$ by \eqref{eq:l3-prime-sw}, \eqref{eq:l3-prime-norm} and
$A_{\mathrm{SW}}=A_0+168$. Interior boxes have constant order predicates,
so $nm=e$ has one labelled representation and $b_m\in\{0,1\}$. Thus
$\|a\|_2\|b\|_2\le(2x_B)^{1/2}$. On each interior box,
$\mathcal S_B(\mathfrak b)=\sum_qC_{\mathfrak b}(q)\Delta_B(q,N)$ with
$\Delta_B=\sum_{\mathrm{pieces}}\Delta_{\mathrm{fs}}$, and
Fact~\ref{fact:bv-full-support} gives
$|\mathcal S_B(\mathfrak b)|/M_c\ll_\sigma NL^{-3692}$.
Fact~\ref{fact:bv-full-support} is the full-support bilinear estimate
derived from Bombieri--Friedlander--Iwaniec, Theorem~0(b). We verified its hypotheses
above for this source.

\subsection{The source-level row-three estimate}\label{subsec:l3-source-row3}

Row three keeps the full anatomy level.
Its factorization keeps every labelled occurrence.
The smooth divisor stays in the first factor.
For $d_*=\rho(\Theta_c-\eta_L)$ and $e_i=\rho t_i$,
the prefix inequalities are
\begin{equation}\label{eq:l3-prefix-margins}
\begin{aligned}
 d_*-2e_1&\ge\frac{33109}{160000}\rho,\\
 d_*-e_1-2e_2&\ge\frac{1859}{160000}\rho,\\
 d_*-e_1-e_2-2e_3&\ge\frac{33109}{480000}\rho
                                      \quad(\chi_c=1).
\end{aligned}
\end{equation}
The anatomy branch gives $\Theta_c\ge(1+t_1)/2$.
The condition $\chi_c=1$ gives the last inequality.
In particular, with $A_3=33109/160000$,
\begin{equation}\label{eq:l3-row3-global-margins}
\begin{gathered}
 \ell_c\ge\max\{A_3-t_3,2\rho t_3\}\ge\frac{33109}{240000},\qquad
 \ell_c-\ell_0\ge\frac{231463}{24720000}=0.0093633899676\ldots,\\
 \Theta_c\le H_{\max}=\frac{153}{256},\qquad
 e_Q=(\rho+\sigma)(153/256-\eta_L)<\frac35,\\
 \frac35-e_Q\ge\frac{391}{160000}
       -\frac{95609}{160000}(\sigma+\tau)>0.0024437499.
\end{gathered}
\end{equation}
Here $e_Q$ bounds the full modulus exponent
$\rho(t_1+t_2+t_3)+(\rho+\sigma)\ell_c$.
The inherited anatomy and rough-factor estimates are Paper~II,
(4.20) and Appendix~B, Lemma~B.1(1)--(2)
\cite[(4.20), Appendix~B, Lemma~B.1(1)--(2)]{PaperII}.
For the tail exponents $s_i$, the cubic conditions give
\begin{equation}\label{eq:l3-tail-margins}
 \sum_{i<l}s_i+2s_l<\ell_c,\qquad
 h=d_*+60\sigma,\qquad
 \Theta_c-h\ge\eta_L-\tau H_{\max}-60\sigma>0.
\end{equation}
Theorem~4.1 of Paper~I then supplies the load surpluses
\begin{equation}\label{eq:l3-packing-loads}
 \left(\frac{7527}{256},\ 60,\ \frac{114421}{1280},\
                         \frac{26247}{256}\right)\sigma.
\end{equation}
The required length data are
\begin{equation}\label{eq:l3-lengths}
\begin{gathered}
 1\le|mn-N|\le N<X^3,\quad
 N_I\ge\frac{T}{6L^{334}}\ge(2T)^{1-2^{-126}},\\
 X^{60\sigma}\le T\le X^{1/3+15\sigma},\quad
 QX^{\sigma/3}\le N_I,\quad QRS<M_IX^{-89\sigma},\quad M_I\ge X^{1/2}.
\end{gathered}
\end{equation}
Here $N_I$ is the selected integer interval length.
The complement has length scale $M_I$.
The letters $Q,R,S$ are the three factor scales.
Each condition holds only above its fixed-$\sigma$ application threshold.
The floor $2\lceil e^{2^{150}}/2\rceil$ is imposed separately.
See Paper~I, (3.6)--(3.9) and (4.9)
\cite[(3.6)--(3.9), Theorem~4.1, (4.9)]{PaperI}.
We verify them for this source below.

The remaining analytic estimate is
\begin{equation}\label{eq:l3-masked-use}
 |\mathcal S_B(\mathfrak b)|
 \ll_{\sigma,A_{\mathrm{loc}}}XL^{-A_{\mathrm{loc}}}
          +X^{1-\ell_0\sigma^2/4}L^{C_{\mathrm{loc}}}
 \quad(\mathfrak b\text{ in row }3).
\end{equation}
It is uniform over the permitted boxes and blocks.
Here $A_{\mathrm{loc}}>0$ is arbitrary;
$C_{\mathrm{loc}}$ is fixed by the local theorem.
Paper~I, Theorem~3.2, (3.13), supplies the local dispersion theorem.
Its Theorem~2.8 enters only in the inner spectral stage.
There we use the form (P1.7), Fact~\ref{fact:spectral-implication}, with the full bracket
\eqref{eq:t28-main-bound}. The inner substitution is
\begin{equation}\label{eq:l3-spectral-dictionary}
 r_T=z,\quad s_T=p_0,\quad n_T=p_0m,\quad
 a_T=\epsilon_yN\quad(\epsilon_y=\pm1),\quad
 A_T=X^{1+\sigma/20}.
\end{equation}
The auxiliary $z$ and $p_0$ are those of Paper~I's dispersion proof.
Both signs and the mixed cusp pair are kept.
Let $\chi$ be smooth, zero below zero and one above one.
For the fixed bump $\psi_0$, define
\begin{equation}\label{eq:l3-nine-weights}
\begin{aligned}
 \psi_{0,1}(v)&=\psi_0(v)\{1-\chi((v-0.8)/0.1)\},\\
 \psi_{0,2}(v)&=\psi_0(v)\chi((v-0.8)/0.1)
                                  \{1-\chi((v-1.4)/0.2)\},\\
 \psi_{0,3}(v)&=\psi_0(v)\chi((v-1.4)/0.2).
\end{aligned}
\end{equation}
They sum to $\psi_0$.
Their strict dyadic windows are
$(0.45,0.90]$, $(0.80,1.60]$, and $(1.40,2.80]$.
The two smooth modulus coordinates use all nine products
$\psi_{0,i}(b/D)\psi_{0,j}(c/C)$, $1\le i,j\le3$.
Multiply each by the fixed cutoffs on the supported $n_T,z,p_0$.
These cutoffs equal one on the actual coefficient support.
The derivative hypothesis is \eqref{eq:t28-seminorms} with
$\epsilon_0=\sigma/2000$ and $\epsilon=\sigma/100$.

Write $\theta_{\mathrm{sp}}=7/32$ and
$V_{\mathrm{sp}}=T^4R^4S^6$.
The coefficient in (P1.7) is zero outside the actual indices
$n_T=p_0m$, $r_T=z$, $s_T=p_0$.
Use its supported norm
\begin{equation}\label{eq:l3-supported-norm}
\begin{gathered}
 \sum_{n''\mid N^\infty}(n'')^{\theta_{\mathrm{sp}}}
                \|\widetilde b(n'')\|_2^2,\\
 \|\widetilde b(n'')\|_2^2
 =\sum_{\substack{n',r,s\\(n',N)=1}}|b_{n'n'',r,s}|^2.
\end{gathered}
\end{equation}
Here $b_{n,r,s}$ is the coefficient of (P1.7).
It is distinct from the complementary sequence $b_m$.
The norm includes no coefficients outside that support.
After substitution, the four bracket contributions divided by
$V_{\mathrm{sp}}$ are bounded by
\begin{equation}\label{eq:l3-four-spectral-terms}
\begin{aligned}
 K_A&=A_T^{\theta_{\mathrm{sp}}}\frac{TQR^2S^5}{X^2},
 &K_2&=\frac{A_T}{X}K_A,\\
 K_R&=\left(\frac{TR}{S}\right)^{2\theta_{\mathrm{sp}}}
                              \frac{TQR^2S^5}{X^2},
 &K_4&=\frac{TQR^2S^5}{X^2}\frac{T^2R^2}{X}.
\end{aligned}
\end{equation}
Only uniform $X^{o(1)}$ factors and powers of $p_0$ are suppressed.
These are precisely Paper~I, (3.81), (3.85), and (3.87)
\cite[(3.69), (3.81)--(3.87)]{PaperI}.
Each term carries the actual-use hypotheses.
\subsection{The row-three local dispersion estimate}\label{subsec:l3-local-row3}

The source-wide estimate also has boundary and comparator contributions.
The local argument gives only its interior row-three block.

We treat one row-three block with $\chi_c=1$ and one full
interior prime-coordinate mesh box. We do not assert the source-wide
boundary, comparator, or row-zero-to-two estimates.
Fix $\sigma>0$, $\tau=\sigma^9$, $\rho=1+\tau$, and a retained
closed clipped parent $c$, used only for the displayed infimum and endpoint
checks. The actual prime cells are half-open. Keep its inherited labels
$t_1\ge t_2\ge t_3$,
$\Theta_c=\inf_c\widetilde H$, $\ell_c=\Theta_c-\eta_L-t_1-t_2-t_3$,
$D_c=N^{\ell_c}$, $z_c=N^{t_3}$ and $u_c=D_c^{\sigma^2}$.
On the whole block $\chi_c=1$. The inequalities
\eqref{eq:l3-prefix-margins}--\eqref{eq:l3-row3-global-margins}
apply with the inherited labels, including clipped endpoints.
Here an interior source box is a full relative-width
$x=(\log N)^{-334}$ mesh box whose closure lies on one side of every
source, order and descriptor face. A clipped terminal mesh fragment
belongs to the boundary contribution.
LEN is not applied to it.

The prefix constants can be checked without an unclipped replacement.
Put $A_3=33109/160000$. Since $t_1,t_2\le25/128$ and
$\Theta_c\ge(1+t_1)/2$, the first two gaps are at least
$\rho(1/2-3(25/128)/2-\eta_L)=\rho A_3$ and
$\rho(1/2-5(25/128)/2-\eta_L)=\rho\,1859/160000$.
For the third, $\ell_c-t_3\ge A_3-2t_3$ and retention gives
$\ell_c-t_3\ge(2\rho-1)t_3$; their maximum is at least $A_3/3$.
Likewise $\ell_c\ge\max\{A_3-t_3,2\rho t_3\}\ge2A_3/3>\ell_0$.
Thus all three prefix WF inequalities have strict room. The whole
anatomy infimum is taken before this branch bound is used.

\paragraph{Exact coefficient and labelled factorization (H1--H2).}
Write $\mathfrak b=(3,c,f,\mathbf D)$. Give the three prefix
occurrences their independent prime sets
$[N^{t_i},N^{\rho t_i})\cap[N^u,N^{\mu_*})\cap\mathcal P_N$.
For $f=2$, append each tail occurrence, independently, with prime set
$[D_j,\min(D_j^\rho,z_c))\cap\mathcal P_N$.
For $f=1$ the tail list is empty. In both cases put $J$ equal to this
complete labelled list, and set $\lambda_d=\lambda_d^-(\mathbf D)$
or $\phi_d^-$, respectively. Thus, exactly,
\[
 C_{\mathfrak b}(q)=\sum_{\substack{d<D_c^\sigma\\d\mid P_N(u_c)}}
 \lambda_d\sum_{(p_j)_{j\in J}}1_{q=d\prod_{j\in J}p_j},
 \qquad |\lambda_d|\le1 .
\]
Neither distinctness mask is inserted here. The orbit divisor $M_c$,
the source weight $w_C$, and the two distinctness corrections remain
outside this coefficient, as in \eqref{eq:l3-masks}--\eqref{eq:l3-adv-split}.
In particular its supports include repeated-prime moduli, and
$(q,2N)=1$. The prime sets are the actual clipped sets, not their
unclipped parents. All labelled primes are at least $u_c\ge u_0$
with $u_0=N^{\ell_0\sigma^2}$, since $\ell_c>\ell_0$ and
$\ell_c\sigma^2<u$.

For packing use the actual upper exponents of the clipped prime sets.
The first three are at most $\rho t_i$; all tail upper exponents are
at most $t_3$. The prefix masks in \eqref{eq:l3-prefix-margins}
and the tail masks $\sum_{i<j}s_i+2s_j<\ell_c$ show that their
merged decreasing list satisfies $\mathrm{WF}(d_*)$, where
$d_*=\rho(\Theta_c-\eta_L)$. For a tail rank, this follows after
adding at most $\rho(t_1+t_2+t_3)$ to its tail inequality.
The top three upper exponents are the prefix endpoints, belonging to
the closed clipped parent. Hence their anatomy value is at least
$\Theta_c$. Put $h=d_*+60\sigma$, $\delta=30\sigma$ and
$g=\log_N T-\delta$.
Every point $\mathbf x$ of the closed clipped parent satisfies
$u\le x_1\le\mu_*=25/128$. On this interval the $v_{\rm pk}(x_1)$ branch of
Paper~I gives $\widetilde H(\mathbf x)\ge v_{\rm pk}(x_1)=(1+x_1)/2$;
hence $\Theta_c\ge(1+u)/2$. Therefore
 $h\ge\Theta_c-\eta_L+60\sigma
 \ge\frac{1+u}{2}-\eta_L>\frac12$; also
 $h<\Theta_c\le153/256$, since
$\Theta_c-h\ge\eta_L-\tau H_{\max}-60\sigma>0$.
Theorem 4.1 of Paper I applies to this same labelled list, with
$\delta\le g\le1/3-\delta/2$. Let its partition be
$J=J_1\sqcup J_2\sqcup J_3$, with the smooth divisor attached to
slot one. Define, using precisely the preceding prime sets,
\[
 \gamma_1(v)=\frac1{|J_1|!}\sum_{d,(p_j)_{J_1}}
             \lambda_d1_{v=d\prod_{J_1}p_j},\qquad
 \gamma_s(v)=\frac1{|J_s|!}\sum_{(p_j)_{J_s}}
             1_{v=\prod_{J_s}p_j}\quad(s=2,3).
\]
Empty slots two and three mean the atom $1_{v=1}$. Slot one always
keeps the $d$ sum, even if $J_1$ is empty. The finite tuple map gives
\[
 C_{\mathfrak b}=C_{\rm slot}(\gamma_1*\gamma_2*\gamma_3),\qquad
 C_{\rm slot}=\prod_s|J_s|!,\qquad |\gamma_s|\le1.
\]
The primes below $u_c$ determine $d$ uniquely. At most $|J_s|!$
assignments have a fixed slot product. The count allows repeated primes.
The rough coordinate count is less than $3/(5\ell_0\sigma^2)$,
so $C_{\rm slot}=O_\sigma(1)$ and
$|C_{\mathfrak b}(q)|/M_c\le M_\sigma!$, with
$M_\sigma=\lceil3/(5\ell_0\sigma^2)\rceil$.
For the tested boxes below, take $C_{\rm ctr}=2$. Uniformly for
sufficiently large $N$,
\[
 \frac{N_I}{TL^{-334}}\in[1/2,2],\qquad
 \frac{M_I}{(X/T)L^{-334}}\in[2,20].
\]

\paragraph{The same tested box (H3--H5).}
Use exactly \eqref{eq:l3-direct-prime} and
\eqref{eq:l3-box-discrepancy}: $X=\prod_iA_i$, $T=\lceil A_1\rceil-1$,
$U_J=\lceil B_1\rceil-1$, and
$\alpha_n=1_{\mathbb P}(n)1_{T<n\le U_J}$.
No product truncation is added to $\Delta_B$. Take $C_{\rm box}=334$;
$\log X=\log N+O(1)$ throughout this family.
For sufficiently large $N$, $T>2^{50}$ and $U_J\le2T$, so
\[
 \alpha_n=1_{P^-(n)\ge T^{51/100}}1_{T<n\le U_J}.
\]
Thus (3.8), option $1$, holds with $\zeta=T^{51/100}\ge z_0(X)$.
The complementary sequence is the actual $b_m\in\{0,1\}$ of
\eqref{eq:l3-direct-prime}; in particular $|b_m|\le\tau_{23}(m)$.
There is no Heath--Brown expansion, logarithm freezing, or proper-power
replacement error for this selected direct-prime representation.
The selected-factor Siegel--Walfisz bound is Paper I (3.25), uniform
in this stronger cutoff and every endpoint and auxiliary coprimality.
Its relative-norm alternative is kept as in (3.48). We assume no lower
bound on any coprimality-truncated selected norm.
In the large-norm case of (3.48), use Lemma~\ref{hyp:ra-ms-transfer}(ii).
Its premises
(3.1), (3.2), (3.6), (3.8) with LEN, (3.11), (3.25), (3.26) and (3.27)
are the data verified in this subsection.

Let $M=\prod_{i=2}^kA_i$, $M_+=\prod_{i=2}^kB_i$,
$M_I=M_+-M$, $m_0=(M+M_+)/2$, and
$\psi_I(m)=\phi_0((m-m_0)/M_I)$, with the fixed bump of Paper I (3.9).
Then $\psi_I=1$ on the full complementary product interval. Since
$5\le k\le11$ and the mesh boxes are full,
$M_I\asymp ML^{-334}$, $m_0\le2X/T$, and
$N_I=U_J-T\ge T/(6L^{334})$, where $L=\log N$.
For fixed $\sigma$ enlarge the threshold so that
$\sigma L\ge2^{130}(\log L+1)$; then
\[
 N_I\ge(2T)^{1-2^{-126}},\quad
 X^{60\sigma}\le T\le X^{1/3+15\sigma},\quad M_I\ge X^{1/2}.
\]
These follow from $N^\alpha/2\le T<N^{1/5}$ and
$\varepsilon N\le X\le N$. The integer convention is
$(T,U_J]$; no continuous interval length replaces $N_I$.

The doubled support lies in
$[M-M_I/2,M_++M_I/2]$. Since $k-1\le10$,
$(1+x)^{k-1}-1\le11x$ eventually. On this entire support and for
every selected integer, $mn\le N_{\rm cut}+25xN$. Consequently
\[
 1\le N-mn,\qquad N-mn\ge(L-25)NL^{-334}>0,\qquad
 |mn-N|\le N<X^3.
\]
This verifies nonresonance on the full doubled rectangle, including
both Cauchy--Schwarz copies. Checking only where $b_m\ne0$
would not suffice.

After localization, $Q,R,S$ denote the scales of
$\gamma_1,\gamma_3,\gamma_2$, respectively. Theorem 4.1 gives, before
fixed dyadic rounding and base conversion, reductions
\[
 (7527/256,\ 60,\ 114421/1280,\ 26247/256)\sigma
\]
in $Q/T$, $T^2R^2S$, $A_0^{7/32}TQR^2S^5$ and
$(TR/S)^{7/16}TQR^2S^5$, respectively, relative to their bounds
$1,N,N^2,N^2$. Here $A_0=X^{1+\sigma/20}$.
The first slot includes the additional load $\sigma\ell_c\le\sigma H_{\max}$.
Each listed reduction strictly exceeds the required losses
$29,8,15,15$. Hence fixed dyadic constants and
$\log_N(N/X)=O(1/L)$ are absorbed for fixed $\sigma$.
In particular the exact requirements (3.6)--(3.7) hold:
\[
 Q<TX^{-29\sigma},\quad N_I\ge QX^{\sigma/3},\quad
 QRS<M_IX^{-89\sigma},\quad M_I\ge X^{1/2},
\]
\[
 T^2R^2S<X^{1-8\sigma},\quad
 A_0^{7/32}TQR^2S^5<X^{2-15\sigma},\quad
 (TR/S)^{7/16}TQR^2S^5<X^{2-15\sigma}.
\]
For $QRS<M_IX^{-89\sigma}$ use the positive fixed gap
$4/5-153/256$ and $M_I\gg N^{4/5}L^{-334}$.
The complete modulus, including $d$ and all labelled primes, obeys
\[
 q<N^{e_Q},\qquad e_Q=(\rho+\sigma)(153/256-\eta_L)<3/5.
\]
Because $X\ge\varepsilon N$, impose
$(3/5-e_Q)\log N\ge(3/5)\log(1/\varepsilon)$ to obtain
$q\le X^{3/5}$. Take the moving frequency $a=N$ and
$K_{\rm scale}=1$: eventually
$0<N\le X\log N\le A_0$ and every supported modulus is prime to $2a$.
This proves H1--H5 with the support relaxation of Paper I Theorem 3.2.

\paragraph{Application and restoration.}
Apply Paper I Theorem 3.2, (3.13), to this full unrestricted coefficient
and this full-support natural-coprime discrepancy. For every fixed
$A_{\rm loc}>0$ it gives exactly
\[
 |\mathcal S_B(\mathfrak b)|
 \ll_{\sigma,A_{\rm loc}}X(\log X)^{-A_{\rm loc}}
       +X^{1-\ell_0\sigma^2/4}(\log X)^{C_{\rm loc}}.
\]
Constants are uniform in these boxes and in blocks with $\chi_c=1$ after fixing
$\sigma$ and source constants. The threshold is ineffective.
Theorem 3.1's smaller abstract support class is not used.
The proof includes all local restorations (3.32)--(3.50):
with $\epsilon=\ell_0\sigma^2/8$ and $A_{\rm work}=A_{\rm loc}+4$,
\[
 XL^{-A_{\rm work}}+XL^Cz_0^{-1/2}+XL^C(Q/N_I)^{1/2}
 +XX^{\epsilon/2}L^C\{u_0^{-1/2}+(Q/N_I)^{1/2}\}.
\]
Keep both terms in braces. They come from restoring the five
cross-copy gcd failures of (3.30), the close-pair set, and the zero-mode
main-term cancellation. The spectral
estimate alone does not give them. The $z_0$ term joins the arbitrary logarithmic saving,
and the remaining terms have savings exceeding $\ell_0\sigma^2/4$.

Theorem 2.8 enters only at Paper I (3.69)--(3.90). In that inner stage,
with $p_0=\min\{p:p\nmid2N\}$, set
$r_T=z$, $s_T=p_0$, $n_T=p_0m$, $a_T=\varepsilon_yN$,
$N_T=8p_0H_*SF$ and $A_T=A_0$. Both signs are kept. The coefficient
is exactly $\beta_{z,m}^{\chi,\varepsilon_y}$ of (3.68) on these
indices and zero elsewhere. In particular
\[
 \sum_{n''\mid N^\infty}(n'')^{7/32}
       \|\widetilde b(n'')\|_2^2
 \ll X^{o(1)}(H_*S)^2F,\quad
 \|\widetilde b(n'')\|_2^2
 =\sum_{\substack{n',r,s\\(n',N)=1}}|b_{n'n'',r,s}|^2,
\]
with no coefficients added outside that support. The unweighted norm
has the same bound. These are the endpoint norm bounds (3.70)--(3.81).
When slot two is a unit, $S=1/2$ and every nonzero-frequency block is
empty. Otherwise (3.82)--(3.85) identify the exact phase, permit
$p_0\mid c$, and use all nine dyadic weights from
\eqref{eq:l3-nine-weights}, with $\epsilon_0=\sigma/2000$ and
$\epsilon=\sigma/100$. The full bracket of
\eqref{eq:t28-main-bound}, including its moving-frequency and mixed-cusp
terms, yields all four quantities in \eqref{eq:l3-four-spectral-terms}.
No bracket term is discarded. The full bounds (3.87)--(3.90) then
control the inner spectral term. After that, the restoration terms
above and $C_{\rm slot}$ are put back.

Finally, this estimate concerns $C_{\mathfrak b}\Delta_B$.
The outer identity is still
\[
 R_{\rm adv}=E_{\rm disp,1}+E_{\rm disp,2}
             +E_{\rm sq}+E_{\rm nc}.
\]
Thus $M_c^{-1}$ and $w_C$ enter once. The forbidden and allowed-repeat
masks remain in $E_{\rm sq}$. The natural-to-$X_A$ comparator remains
in $E_{\rm nc}$. The source cutoff, peel, boundary and terminal-mesh
restorations, capped rows zero to two, and complete-family coefficient
counts require their own estimates. This local estimate alone
gives neither the complete L3 bound nor the RSW bound.

The outer application of Theorem~3.2 above supplies \eqref{eq:l3-masked-use}.
The scoped spectral implications in Section~\ref{sec:classical-inputs}
do not by themselves give it.

\subsection{The complete box family}

The required complete-family bounds are
\begin{equation}\label{eq:l3-boundary}
 \#\mathcal E_{\mathrm{bdry}}\ll NL^{-334},\qquad
 D_{\mathrm{bdry}}\ll_\sigma NL^{-331/2},\qquad
 \sum_B|w_B|\ll_\sigma L^{3689}.
\end{equation}
There are at most $122$ face families per $k$.
Their integer normals have maximum coefficient five.
Their coefficient sums are at most eleven.
\paragraph{Proof of the complete-family bounds.}
Use the unclipped geometric mesh $[A^{(j)},A^{(j+1)})$,
$A^{(j)}=(N^\alpha/2)(1+L^{-334})^j$, $j\ge0$, in every coordinate.
Boxes with a coordinate $\ge N^{1-4\alpha}$ carry no source mass, since
$k\ge5$ and every prime is at least $N^\alpha$. So there are no terminal
fragments, and every interior box is a full mesh box.
\emph{Box count.} Each coordinate has at most
$((1-\alpha)L+1)(L^{334}+1)+1\le L^{335}$ mesh intervals for $L\ge24$.
With $k\le11$, finitely many cells and $w_C\le36$,
$\sum_B|w_B|\le36|\mathcal C_+|\sum_{k=5}^{11}L^{335k}\ll L^{3685}\le L^{3689}$.
\emph{Boundary count.} Let $h'=\log(1+L^{-334})/L\le L^{-335}$.
A source integer $e$ in a boundary box has its vector within $h'$
in each coordinate of a face. A radial face forces $e$ into
$[y(1+L^{-334})^{-11},y(1+L^{-334})^{11}]$ with $y\in\{\varepsilon N,N_{\mathrm{cut}}\}$.
There are at most $2(23NL^{-334}+1)$ such $e$.
Every other face is $c_ix_i+c_jx_j=\theta$ with integers $|c_i|,|c_j|\le5$,
$c_i\ne0$ \cite[(5.29a)--(5.31) and the facet discussion after (5.40)]{PaperII}.
Write $e=p_im'$. For fixed $m'$, $p_i$ lies in an interval of relative
width at most $21L^{-334}$ below $N/m'$, which contains at most
$21L^{-334}N/m'+1$ integers. Here $m'$ is a product of at most ten primes
in $[N^\alpha,N)$, and $m'<N^{1-\alpha}$. By Mertens,
$\sum_{m'}1/m'\le\exp(\sum_{N^\alpha\le p<N}1/p)\le e/\alpha<32$ eventually.
Hence one face family contributes at most $700NL^{-334}+N^{1-\alpha}$.
The number of face families is finite ($\le122$ for each $k$). Thus
$\#\mathcal E_{\mathrm{bdry}}\le C_{\mathrm{bd}}NL^{-334}$ with an absolute
$C_{\mathrm{bd}}$. No numerical value is required below.
\emph{Boundary discrepancy.} Let $A_{\mathrm{bd}}$ be $A$ restricted to
boundary boxes. Then
$D_{\mathrm{bdry}}=\sum_{\mathfrak b}M_c^{-1}\langle C_{\mathfrak b},\Delta_{P,\mathrm{bd}}\rangle$.
By $|C_{\mathfrak b}|/M_c\le C_{\mathrm{L3},\sigma}$,
$\sum_q1_{q\mid N-e}\le\tau(N-e)$, Cauchy--Schwarz,
$\tau^2\le\tau_4$, $\sum_{n\le N}\tau_4(n)\le N(1+L)^3$ and
$\sum_{q\le N}1/\varphi(q)\le2(1+L)$,
\[
 |D_{\mathrm{bdry}}|\ll_\sigma36(\#\mathcal E_{\mathrm{bdry}})^{1/2}(N(1+L)^3)^{1/2}
             +36\,\#\mathcal E_{\mathrm{bdry}}(1+L)\ll_\sigma NL^{-331/2},
\]
with finitely many blocks at fixed $\sigma$. This proves
\eqref{eq:l3-boundary}, with the first bound in the form
$\#\mathcal E_{\mathrm{bdry}}\ll NL^{-334}$.
Under \eqref{eq:l3-masked-use}, take $A_{\mathrm{loc}}\ge3692$.
The $L^{3689}$ box count and $W\gg N/L^2$ then leave $L^{-1}$.
Equations \eqref{eq:l3-bv-total} and \eqref{eq:l3-boundary} then give
\begin{equation}\label{eq:l3-disp-total}
 (R_{\mathrm{adv}})_+/W=o_\sigma(1),
\end{equation}
provided \eqref{eq:l3-repeated-cap} and the coprime-comparator bound hold.

\subsection{The principal term and its price}

\paragraph{(P2.8d) The source-independent profile comparison.}
We use Paper~II, (4.70)--(4.75), in the following form.
Put
\begin{equation}\label{eq:l3-profile}
\begin{aligned}
 v_1(s)&=\min\left\{\frac{1+s}{2},\frac{32-34s}{39},
                                      \frac{57}{64}-\frac32s\right\},\\
 v_*(s)&=v_1(s)-u,\qquad
 \psi_*(x,y)=\begin{cases}x,&x\ge y,\\-\infty,&x<y,\end{cases}\\
 \Theta_*(x,y,z)&=\max\{v_*(x),v_*(y),v_*(x+y),v_*(x+y+z),\\
 &\hspace{14mm}\psi_*(v_*(x+z),x+2y+z),
                    \psi_*(v_*(y+z),2x+y+z)\}.
\end{aligned}
\end{equation}
The subtraction of $u$ precedes both guards.
This is the same guard convention as
\hyperref[eq:app-six-branch]{equation~\ref*{eq:app-six-branch}}.
It agrees with Paper~II's zero guard because an unguarded branch is
positive, as explained immediately after that display.
Define
\begin{equation}\label{eq:l3-Gstar}
\begin{gathered}
 H_{\log}=\log(\mu_*/u),\qquad
 a_3(\mathbf t)=\frac{\Theta_*(\mathbf t)-t_1-t_2-t_3}{t_3},\\
 J_*=\int_u^{\mu_*}\frac{dt_1}{t_1}\int_u^{t_1}\frac{dt_2}{t_2}
                 \int_u^{t_2}\frac{1-f(a_3(\mathbf t))}{t_3^2}\,dt_3,\\
 R_{\mathrm{up}}=500(F(30)-1)+500H_{\log}(1-f(30))
                         +250H_{\log}^2(F(30)-1),\\
 G_*=\frac1{\mu_*}+R_{\mathrm{up}}+J_*.
\end{gathered}
\end{equation}
Extend $f$ by zero below two.
These are Paper~II, (4.70)--(4.75)
\cite[(4.70)--(4.75)]{PaperII}.
Their source-independent comparison uses
\begin{equation}\label{eq:l3-profile-margin}
 g=(u-\eta_L)-\frac92\mu_*\tau>0.001899>2\mu_*\tau.
\end{equation}
Hence a positive canonical lower term lies in a cell with $\chi_c=1$.
The capped rows still satisfy $s_c>30$.
Their residuals are bounded by the same $R_{\mathrm{up}}$.
Thus the cap changes neither $G_*$ nor its price.

Let $\mathcal P_\Pi$ be the signed prefix profile:
\begin{equation}\label{eq:l3-principal-profile}
\begin{aligned}
 \mathcal P_\Pi={}&\frac1uF(s_0)
 -\frac1u\sum_{c:j=1}\operatorname{vol}_c f(s_c)
 +\frac1u\sum_{c:j=2}\operatorname{vol}_c F(s_c)\\
 &-\sum_{c:j=3,\chi_c=1}\frac{\operatorname{vol}_c}{t_3}f(\ell_c/t_3),
 \qquad \operatorname{vol}_c=\int_c\prod_i\frac{dt_i}{t_i}.
\end{aligned}
\end{equation}
The required mass and comparison statements are
\begin{equation}\label{eq:l3-principal-estimates}
\begin{gathered}
 X_A\le\frac NL\{(1-\varepsilon)I_{\mathrm{HL}}+O(L^{-1})\},\\
 M_{\mathrm{prin}}=X_A\frac{2e^{-\gamma}C(N)}L
                            (\mathcal P_\Pi+o_\sigma(1)),\qquad
 \mathcal P_\Pi\le G_*+o_\sigma(1).
\end{gathered}
\end{equation}
The first bound is Paper~II, (5.34)--(5.42).
The last comparison uses (4.72)--(4.75)
\cite[(4.72)--(4.75), (5.34)--(5.42)]{PaperII}.
\paragraph{Proof of \eqref{eq:l3-principal-estimates}.}
\emph{Mass.} The cells $\mathcal C_+$ are Paper~II's cells with $w_C\ge1$,
and the radial range $\varepsilon N<e<N$ is Paper~II's. The outer-box
majorant $f^+$ of Paper~II, (5.34), counts labelled factorizations into
all primes in $[N^\alpha,N)$, and $c(\mathbf j)\ge w_C$ for the box
$\mathbf j$ containing the ordered factor vector of $e$, which lies in
$C_N$. Hence $A\le f^+$ pointwise directly (no comparison with the
intermediate $g$ is used). Paper~II, (5.34)--(5.42), then gives
$X_A\le X^+=\frac NL\{(1-\varepsilon)I_{\mathrm{HL}}+O(L^{-1})\}$
\cite[(5.32), (5.34), (5.42)]{PaperII}. In particular $X_A\le N/L$
eventually, since $I_{\mathrm{HL}}\le\overline H<1$.
\emph{Principal term.} As in Paper~II, Lemma~B.9,
$V_N(z)=(1+o(1))2e^{-\gamma}C(N)/\log z$ uniformly for $z\ge N^h$,
$h=1/1000$ \cite[(B.9) and its proof]{PaperII}. The proof of (B.9) gives
an exact factorization. Mertens' theorem gives the asymptotic. The
primes $p\mid N$ with $p\ge z$ change the product by $1+O(\omega(N)/z)$.
Every actual cutoff is at least $N^{u/\rho^3}>N^h$ \cite[Lemma~B.8]{PaperII}. At fixed $\sigma$ the cells are
finitely many, with fixed exponent boxes. Mertens gives
$\sum_{v\in c}1/\varphi(v)=\operatorname{vol}_c+o(1)$, the primes
dividing $2N$ costing $O(\omega(N)N^{-u})$. Also $K_j(s_c)$ is bounded.
Summing $\nu_jM_j$ gives
$M_{\mathrm{prin}}=X_A\frac{2e^{-\gamma}C(N)}L(\mathcal P_\Pi+o_\sigma(1))$,
with $\log_Nz_{3,c}=t_3$ supplying the factor $1/t_3$.
\emph{Profile comparison.} Write $F=1+(F-1)$, $f=1-(1-f)$.
On rows $0,1,2$, $s_c>30$, and $F$ decreases and $f$ increases. With
$\sum_{j=1}\operatorname{vol}_c=H_{\log}$ and
$\sum_{j=2}\operatorname{vol}_c=H_{\log}^2/2$, the residuals are at most
$R_{\mathrm{up}}$ \cite[(4.73)]{PaperII}. For row three, if
$f(a_3(\mathbf x))>0$ at an actual point $\mathbf x$ of a cell, then
\eqref{eq:l3-profile-margin} gives $\chi_c=1$ and
$\ell_c\ge\Theta_*(\mathbf x)-x_1-x_2-x_3$ \cite[(4.72) and the two
paragraphs after it]{PaperII}. For $x_3\in[t_3,\rho t_3]$, monotonicity of
$sf(s)$ at fixed numerator $\ell_c$ gives
$f(\ell_c/t_3)/t_3\ge f(\ell_c/x_3)/x_3\ge f(a_3(\mathbf x))/x_3$.
Integrating over each retained cell,
$\sum_{j=3,\chi_c=1}\operatorname{vol}_cf(\ell_c/t_3)/t_3
 \ge\int\!\!\int\!\!\int f(a_3)\,dt_1dt_2dt_3/(t_1t_2t_3^2)$.
The constant parts give $1/\mu_*$ by \cite[(4.74)]{PaperII}. Hence
$\mathcal P_\Pi\le1/\mu_*+R_{\mathrm{up}}+J_*+o_\sigma(1)=G_*+o_\sigma(1)$,
which is \cite[(4.75)]{PaperII} for this source.

Use the same coefficient bound as the exterior comparison:
\begin{equation}\label{eq:l3-certificate}
\begin{gathered}
 G_*\le G_0=\frac{6088440363}{10^9},\qquad
 2e^{-\gamma}G_0\le c_{K2,G0}
 =\frac{6836825163875517756286846321650754372931094689}{10^{45}},\\
 I_{\mathrm{HL}}\le\overline H=\frac{389253}{5000000}.
\end{gathered}
\end{equation}
The functional bound \(G_*\le G_0\) is Paper~II,
\cite[(4.73), (4.75), (4.79), (A.10a)]{PaperII}, for the identical
canonical functional \eqref{eq:l3-Gstar}. Its scalar rounding is the
same as \eqref{eq:k2-certificate}. The mass bound is Paper~II, (A.11)
\cite[(A.11)]{PaperII}.
The optional unshifted certificate in
\hyperref[subsec:app-g24]{the unshifted-comparison subsection of
Appendix~B} is outside this proof.

The constant $c_{\mathrm{O1}}$ is optional, non-certified appendix
material outside the L3 proof.
In the optional scalar comparisons below, $c$ denotes the rational
from \eqref{eq:intro-exterior-constant}.
No functional comparison with $c_{\mathrm{O1}}$ is needed here.
Its rational definition is
\begin{equation}\label{eq:l3-exact-certificate-constant}
 c_{\mathrm{O1}}=\frac{\sum_{i=0}^7n_iB^i}{\sum_{i=0}^7d_iB^i},
 \qquad B=10^{25},
\end{equation}
with the integers in Table~\ref{tab:l3-certificate-integers}.
\begin{table}[ht]
\caption{Exact certificate constant}\label{tab:l3-certificate-integers}
\small
\begin{tabular}{@{}rrr@{}}
\hline
$i$ & $n_i$ & $d_i$\\
\hline
 0 & 4790995369589327356095267 & 5033118760359274030825472\\
 1 & 6963256824825772577191322 & 8266496527677353149933429\\
 2 & 8996423827217290029713216 & 9138360880901719297497815\\
 3 & 4916164251930472277375479 & 4671424972330288696699300\\
 4 & 329572634628519073712385 & 939590108467331323910141\\
 5 & 9317769528669888650546075 & 5634060699913223990647374\\
 6 & 8280830054071768648850166 & 2330146450712536415611653\\
 7 & 389762593301325772 & 57030507720277440\\
\hline
\end{tabular}
\end{table}
The exact rational comparison of $c$ and $c_{\mathrm{O1}}$ is
in Section~\ref{sec:final-bill}, \eqref{eq:bill-coefficient-cost}. It uses no exponential
comparison.
It is independent of the coarse enclosure for $U_0$.
The mass-integral bound is Paper~II, (A.11)
\cite[(A.11)]{PaperII}.
Put
\begin{equation}\label{eq:l3-price}
\begin{gathered}
 m_{\mathrm{sw}}=\frac{e^{-\gamma}}2I_{\mathrm{HL}}G_*,\qquad
 \Delta_{\mathrm{cert}}=\frac{c_{K2,G0}}4\overline H-m_{\mathrm{sw}}\ge0,\\
 C_{\mathrm{sw}}=\frac{650000000}{1111}C_E,\qquad
 P_{\mathrm{sw}}=\frac{48544888063158}{27775},\\
 a_{\mathrm{sw}}=C_{\mathrm{sw}}\sigma,\qquad
 p_{\mathrm{sw}}=P_{\mathrm{sw}}\sigma^9.
\end{gathered}
\end{equation}
The remaining required estimates are
\begin{equation}\label{eq:l3-small-estimates}
\begin{gathered}
 (M_{\mathrm{prin}}/W-m_{\mathrm{sw}})_+=o_\sigma(1),\\
 M_{\mathrm{err}}/W\le C_{\mathrm{sw}}\sigma+o_\sigma(1),\qquad
 P_3/W\le P_{\mathrm{sw}}\sigma^9+o_\sigma(1),\\
 |E_{\mathrm{nc}}|/W=O_\sigma(N^{-\alpha}L^4).
\end{gathered}
\end{equation}
The exact numerical multipliers are
\begin{equation}\label{eq:l3-multipliers}
 \frac{4000\cdot585}{4(1-10^{-4})}=\frac{650000000}{1111},\qquad
 \frac{36\cdot502002(459/100)^3}{1-10^{-4}}
                       =\frac{48544888063158}{27775}.
\end{equation}
\paragraph{(P2.8c) Density product and prefix sums.}
Put $h=1/1000$ and $H_N=\sum_{N^h\le p\le N}(p-1)^{-1}$.
The two bounds used in the first multiplier are
\begin{equation}\label{eq:l3-mass-prefix-sums}
\begin{gathered}
 V_N(z)\le\frac{4C(N)}{\log z}\le\frac{4000C(N)}L
                                  \quad(z\ge N^h),\\
 \sum_{\text{actual prefixes }v}\frac1{\varphi(v)}
       \le\sum_{j=0}^3H_N^j\le1+8+64+512=585.
\end{gathered}
\end{equation}
These are Paper~II, Appendix~B, Lemmas B.9 and B.11,
(B.9) and (B.11) \cite[Appendix~B, Lemmas B.9, B.11, (B.9), (B.11)]{PaperII}.
Each half-open prefix cell is counted once.
Using $X_A\le N/L$ from (P2.8d) now gives
\begin{equation}\label{eq:l3-secondary-mass-step}
 M_{\mathrm{err}}\le\frac{X_A4000C(N)}L\,585
               \{C_E\sigma+o_\sigma(1)\}.
\end{equation}
Division by $W=4(1-\varepsilon)C(N)N/L^2$ gives $C_{\mathrm{sw}}$.

\paragraph{(P2.8e) The upper mass and the peel.}
The source mass is \eqref{eq:l3-principal-estimates},
Paper~II, (5.34)--(5.42).
The native lower-sieve cutoff uses the finite peel
\eqref{eq:l3-peel} and the estimate
\begin{equation}\label{eq:l3-peel-restatement}
 \frac{P_3}{W}\le
 \frac{36\cdot502002(459/100)^3}{1-\varepsilon}\tau+o_\sigma(1).
\end{equation}
For the uniform two-form calculation, put $m=vr$ and $w=\min(z_c,N^\alpha)$.
The required ranges are
\begin{equation}\label{eq:l3-selberg-ranges}
\begin{gathered}
 \log_Nm\le153/256+2\tau,\qquad \log_Nw\ge2/1001,\\
 w^2m\le N^{253/256+(153/64)\tau}<N^{1-0.0116}.
\end{gathered}
\end{equation}
For each sieve prime $p<w$, every prime factor of $m=vr$ is at least
$z_c\ge w$, so $p\nmid m$. The local two-form root count is then
$\nu_N(p)=1$ for $p\mid N$, and $\nu_N(p)=2$ otherwise.
For squarefree $d$, define $g_N(d)$ multiplicatively by
$g_N(p)=\nu_N(p)/(p-\nu_N(p))$.
The denominator and one-layer bound are
\begin{equation}\label{eq:l3-selberg-layer}
\begin{gathered}
 \mathcal G_N(w)=\sum_{\substack{d\le w\\d\mid P(w)}}g_N(d)
             =\frac{\log^2w}{4C(N)}(1+o(1)),\\
 S^{(1)}_m(w)\le\frac{N/m}{\mathcal G_N(w)}+O(w^2\log^5w)
       \le502002\,\frac{2C(N)N}{mL^2}.
\end{gathered}
\end{equation}
Here $P(w)$ is the product of all primes below $w$.
The last step uses
$2.001/(2/1001)^2=501251.00025<502002$.
The exponent saving in \eqref{eq:l3-selberg-ranges} absorbs its remainder.
There are at most $36$ source layers.
For the parent-bin prime sums $M_h$, the needed estimates are
\begin{equation}\label{eq:l3-peel-harmonic-sums}
 M_h=\sum_{p\ \mathrm{in\ bin}\ h}\frac1p\le2\tau,\qquad
 \sum_hM_h\le\frac{459}{100},\qquad
 \sum_{a_1,a_2,a_3,r}\frac1{a_1a_2a_3r}
       \le2\left(\frac{459}{100}\right)^3\tau.
\end{equation}
The pair $a_3,r$ belongs to one parent bin.
Multiplication of \eqref{eq:l3-selberg-layer} by $36$ and this sum,
then division by $W$, gives \eqref{eq:l3-peel-restatement}.
These are the exact steps of Paper~II, (4.45)--(4.56),
with the layer ceiling changed from three to $36$
\cite[(4.45)--(4.56), (5.34)--(5.42)]{PaperII}.
The uniform two-form estimates are proved here under the stated ranges.
\paragraph{Proof of \eqref{eq:l3-small-estimates}.}
\emph{Principal.} Since $\mathcal P_\Pi\le G_*+o_\sigma(1)$ and $G_*>0$,
the mass bound gives
$M_{\mathrm{prin}}/W\le\frac{(1-\varepsilon)I_{\mathrm{HL}}2e^{-\gamma}(G_*+o(1))}{4(1-\varepsilon)}
=m_{\mathrm{sw}}+o_\sigma(1)$.
\emph{Secondary sieve.} $E_{s,c}\le C_E\sigma+o_\sigma(1)$ uniformly,
since $\ell_c\ge\ell_1$. Lemma~B.9, Lemma~B.11 and $X_A\le N/L$ give
\eqref{eq:l3-secondary-mass-step}. Every prefix lies in one row, so the
four depths together contribute at most $585$ \cite[(B.9), (B.11)]{PaperII}.
\emph{Peel.} For $m=vr$ in $P_3$, all primes of $m$ are at least $z_c$,
$(m,N)=1$ and $z_c\in[N^{u/\rho^3},N^{\rho\mu_*}]$. Since $\chi_c=1$,
$\log_Nm\le\rho(t_1+t_2+2t_3)\le\rho(153/256-\eta_L)<153/256+2\tau$
\cite[(4.55)]{PaperII}. Write $A=\sum_{h=1}^{36}1_{\{A\ge h\}}$.
Each layer consists of $N^\alpha$-rough $e$. With $N-e=mk$, both $k$ and
$N-mk$ are $w$-rough, $w=\min(z_c,N^\alpha)$. A prime $p<w$ dividing $k$
and $N$ would divide $e$, and $k$ is odd. The two-form Selberg bound
\cite[(4.48)--(4.54)]{PaperII} gives \eqref{eq:l3-selberg-layer} for each
layer, with remainder absorbed by \eqref{eq:l3-selberg-ranges}. The prime
sums \eqref{eq:l3-peel-harmonic-sums} are \cite[(4.56)]{PaperII}, since
$r$ and $a_3$ lie in one parent bin. Multiplying by $36$ gives
\eqref{eq:l3-peel-restatement}.
\emph{Comparator.} For squarefree $q$,
$|d_A(q)|=\varphi(q)^{-1}\sum_{(e,q)>1}A(e)$. A common prime $\ell$ of
$e$ and $q$ is at least $N^\alpha$, so with $q=\ell r$,
\[
 \sum_{q\ \mathrm{sf}}|O_{\mathfrak b}(q)d_A(q)|
 \le C_{\mathrm{L3},\sigma}\sum_{\ell\ge N^\alpha}\frac{36N}{\ell(\ell-1)}
     \sum_{r\le N}\frac1{\varphi(r)}\ll_\sigma N^{1-\alpha}L.
\]
Summing finitely many blocks and dividing by $W$ gives
$|E_{\mathrm{nc}}|/W=O_\sigma(N^{-\alpha}L^3)$, which implies the stated
$O_\sigma(N^{-\alpha}L^4)$. This is the coprime-comparator bound used in
\eqref{eq:l3-disp-total}.

\subsection{The finite signed identity and its consequence}

\paragraph{(P2.8a) Finite signed identities.}
For the finite source \eqref{eq:l3-source} and stated sieve data,
combine \eqref{eq:l3-source-id}, \eqref{eq:l3-buchstab-peel}
and \eqref{eq:l3-aa-id}.
Their conditional signed identity is
\begin{equation}\label{eq:l3-rsw-id}
 S_{\mathrm{sw}}=E_{\mathrm{src}}+P_3+M_{\mathrm{prin}}
   +M_{\mathrm{err}}+R_{\mathrm{adv}}
                   -(Z_{\mathrm{src}}+Z_{\mathrm{AA}}+Z_{\mathrm{peel}}).
\end{equation}
Define
\begin{equation}\label{eq:l3-remainder}
\begin{aligned}
 \delta_P&=P_3/W-p_{\mathrm{sw}},\qquad
 \delta_A=M_{\mathrm{err}}/W-a_{\mathrm{sw}},\qquad
 \delta_M=M_{\mathrm{prin}}/W-m_{\mathrm{sw}},\\
 R_{\mathrm{sw}}&=E_{\mathrm{src}}/W+\delta_P+\delta_A+\delta_M
                                                +R_{\mathrm{adv}}/W,\\
 V_{\mathrm{sw}}&=(R_{\mathrm{sw}})_+\le E_{\mathrm{src}}/W
     +(\delta_P)_++(\delta_A)_++(\delta_M)_++(R_{\mathrm{adv}})_+/W.
\end{aligned}
\end{equation}
Under the stated sieve hypotheses, all three $Z$ terms are nonnegative.
Consequently
\begin{equation}\label{eq:l3-final}
\begin{aligned}
 -\frac{S_{\mathrm{sw}}}{W}
 &=-m_{\mathrm{sw}}-a_{\mathrm{sw}}-p_{\mathrm{sw}}-R_{\mathrm{sw}}
       +\frac{Z_{\mathrm{src}}+Z_{\mathrm{AA}}+Z_{\mathrm{peel}}}{W}\\
 &\ge-\frac{c_{K2,G0}}4\overline H
          -C_{\mathrm{sw}}\sigma-P_{\mathrm{sw}}\sigma^9-V_{\mathrm{sw}}.
\end{aligned}
\end{equation}
The favorable $Z_{\mathrm{peel}}$ is used once.
Every adverse remainder remains in \eqref{eq:l3-remainder}.

Assume all estimates stated above for the same $A$.
Then \eqref{eq:l3-source-error}, \eqref{eq:l3-disp-total}
and \eqref{eq:l3-small-estimates} imply
\begin{equation}\label{eq:l3-remainder-limit}
 V_{\mathrm{sw}}=o_\sigma(1).
\end{equation}
\paragraph{Proof of \eqref{eq:l3-remainder-limit}.}
All estimates above are proved for the same $A$, the same cells and the
same blocks. By \eqref{eq:l3-source-error}, $E_{\mathrm{src}}/W\ll L^{-331}$.
By \eqref{eq:l3-small-estimates}, $(\delta_M)_+,(\delta_A)_+,(\delta_P)_+$
are $o_\sigma(1)$. For $R_{\mathrm{adv}}$, \eqref{eq:l3-adv-split} and
\eqref{eq:l3-interior-split} give
\[
 \frac{|R_{\mathrm{adv}}|}W\le
 \frac{|D^{(0)}_{\mathrm{int}}|+|D^{(1)}_{\mathrm{int}}|+|D^{(2)}_{\mathrm{int}}|
      +|D^{(3)}_{\mathrm{int}}|+|D_{\mathrm{bdry}}|+|E_{\mathrm{sq}}|+|E_{\mathrm{nc}}|}W .
\]
The terms are $O_\sigma(L^{-1})$ by \eqref{eq:l3-bv-total};
$O_\sigma(L^{-1})$ by \eqref{eq:l3-masked-use} with $A_{\mathrm{loc}}=3692$,
$|w_C|/M_c\le36$ and $\sum_B|w_B|\le L^{3689}$;
$O_\sigma(L^{-327/2})$ by \eqref{eq:l3-boundary};
and $O_\sigma(N^{-\ell_1\sigma^2}L^4)$, $O_\sigma(N^{-\alpha}L^4)$.
Hence $V_{\mathrm{sw}}=o_\sigma(1)$.
After fixing $\sigma_*$, choose the eventual threshold so that
\begin{equation}\label{eq:l3-remainder-share}
 V_{\mathrm{sw}}\le\frac1{200000}.
\end{equation}
That threshold may be ineffective.
It also covers
\begin{equation}\label{eq:l3-structural-thresholds}
\begin{gathered}
 N_{\mathrm{dom}}=2\left\lceil\frac12
  \exp\!\left(\max\{10^4,2\log2/(\ell_1\sigma_*^2)\}\right)\right\rceil,\\
 N_{z,0}=2\left\lceil\frac12
       \exp\!\left(2\exp((\ell_1\sigma_*^2)^{-1/3})\right)\right\rceil,
 \qquad N_{\mathrm{LEN,app}}(\sigma_*).
\end{gathered}
\end{equation}
These structural thresholds alone do not imply \eqref{eq:l3-remainder-share}.

Relative to the comparison $-S_{\mathrm{sw}}/W\ge-2\overline H$,
the constant difference is
\begin{equation}\label{eq:l3-gain}
\begin{aligned}
 \Delta D_{L3}&=\left(2-\frac{c_{K2,G0}}4\right)\overline H\\
 &=0.022638464724298154340603\ldots.
\end{aligned}
\end{equation}
The exact cost of using the common coefficient instead of the optional
coefficient is
\[
 \frac{\overline H}{4}(c_{K2,G0}-c)
 =0.000049489563023753677508253183146688036059673516909.
\]
This is scalar arithmetic. Its use in the final count requires \eqref{eq:l3-final}
and the common error bounds.
The capped-row application is \eqref{eq:l3-bv-application}--\eqref{eq:l3-bv-total}.
The row-three estimate is \eqref{eq:l3-masked-use}, proved in the
row-three subsection. All errors use the common source
\eqref{eq:l3-source}. The inputs are Paper~I, Theorem~2.8;
Fact~\ref{fact:bv-full-support}; Lemma~\ref{hyp:ra-ms-transfer}(ii) for
(3.48); and the prime estimates of Fact~\ref{fact:prime-estimates}.
Hence \eqref{eq:l3-final} holds at the
threshold of \eqref{eq:l3-remainder-share}. The constant \eqref{eq:l3-gain} is
a gross coefficient difference. It enters the final count only together
with the switched error $C_{\mathrm{sw}}\sigma_*+P_{\mathrm{sw}}\sigma_*^9
+V_{\mathrm{sw}}$ of \eqref{eq:l3-final}. In the final count
this error is bounded by the caps \eqref{eq:bill-sigma}
($C_{\mathrm{sw}}\sigma_*\le10^{-5}$, $P_{\mathrm{sw}}\sigma_*^9\le10^{-5}$)
and by \eqref{eq:l3-remainder-share}, in total at most $1/40000$.
